\documentclass[conference]{IEEEtran}
\usepackage[T1]{fontenc}
\usepackage{amsmath,amssymb,booktabs,url}
\usepackage{array}
\usepackage{microtype}
\title{Forecastability Is Conditional: A Horizon-Conditioned and Calibration-Aware Study Across Complex Systems}
\author{\IEEEauthorblockN{Anonymous Proposal Author}\IEEEauthorblockA{Design-only research plan generated through the Survey--Idea--ExperimentDesign--Author workflow}}
\begin{document}
\maketitle

\begin{abstract}
Can an advanced intelligent machine predict the future? The answer depends on what ``the future'' denotes, what information is available, and how far ahead the forecast is issued. This paper proposes a Horizon-Conditioned Forecastability and Calibration Atlas (HCFCA), a reproducible benchmark for separating deterministic state prediction, probabilistic forecasting, and conditional scenario projection. The design compares mechanistic, statistical, machine-learning, and hybrid forecasters on orbital state, weather, climate hindcast, economic, and aggregate public-health time series. It normalizes lead time by a domain-specific characteristic timescale, perturbs the information set, evaluates both skill and calibration, and estimates an effective horizon at which a model no longer adds reliable information over a strong baseline. The central claim is deliberately ambitious but testable: intelligent models can extend the useful horizon in observation-rich and structurally stable tasks, yet they cannot remove the horizon imposed by initial-condition uncertainty, model approximation, stochastic forcing, unobserved interventions, and distribution shift. This artifact is a design-only proposal. It reports no executed experiment and contains no observed results.
\end{abstract}

\begin{IEEEkeywords}
forecastability, uncertainty quantification, probabilistic forecasting, complex systems, machine learning, calibration, chaos
\end{IEEEkeywords}

\section{Introduction}
The popular question ``Is it possible to predict the future?'' contains several scientific questions that should not be collapsed into one. The position of a satellite can often be propagated from a well-measured state under a validated dynamical model. A weather service instead issues a distribution over possible atmospheric states because small analysis errors and model approximations amplify with lead time. A climate projection describes conditional distributions under specified forcings, while an economic or public-health forecast can change the system it attempts to describe because people, institutions, and policies respond to information.

The distinction is consequential. A point estimate can be accurate on average while giving unreliable probabilities. A long-range scenario can be useful without being a prediction of an exact event. A fast machine-learning model can process more forecasts without acquiring information that was never observed. The scientific objective is therefore not to build a universal oracle, but to measure when predictive information remains, how it decays, and whether an intelligent machine can recognize the boundary of its own useful forecast.

Lorenz demonstrated that finite deterministic nonlinear systems can generate nonperiodic flow and sensitive dependence on initial conditions \cite{lorenz1963,lorenz1969}. In operational weather prediction, initial-condition uncertainty and model approximation are recognized sources of lead-time error growth, motivating ensemble prediction \cite{ecmwf,leutbecher2008}. These results do not imply that all futures are random. They imply that predictability is conditional on the system, target, information set, and horizon.

Modern machine learning changes the engineering frontier. GraphCast reported global medium-range forecasts that were more accurate than a leading operational deterministic system on most of its stated verification targets and could be generated rapidly \cite{lam2023}. That is a real advance for a bounded task. It is not evidence that computation abolishes chaotic amplification, nonstationarity, or unknown interventions. The present proposal turns this boundary into the primary object of study.

\section{Background, Survey, and Research Gap}
\subsection{Predictability of dynamical systems}
For a state $x_t$ and an evolution operator $\Phi_h$, deterministic prediction has the form $\hat{x}_{t+h}=\Phi_h(x_t)$. If the observed initial state is $x_t+\eta_t$, the prediction error can grow approximately as $\|D\Phi_h\eta_t\|$ in a locally linear regime. Positive local expansion rates make a small measurement error consequential at longer lead times. In a many-scale flow, unresolved scales can feed back into the resolved state, making the error-growth curve depend on both the observation and model resolution.

The relevant limit is not only a property of the equations. It is a property of the pair $(X_t,Y_{t+h})$, where $X_t$ is the available information and $Y_{t+h}$ is the target. A target with strong persistence or a known conservation law can retain information for a long time. A target driven by unobserved impulses can lose it quickly. The same system can be predictable for one variable and effectively unpredictable for another.

\subsection{Forecasts, probabilities, and projections}
ECMWF describes forecast skill as decreasing with lead time because of uncertainty in the initial conditions and approximations required by numerical models; its ensemble system uses perturbed initial conditions and stochastic model tendencies to represent these uncertainties \cite{ecmwf}. The ensemble is not a guarantee that the truth lies inside a stated interval. Its calibration must be checked against held-out outcomes.

The IPCC uses a related but distinct language for climate information. Regional information is distilled from observations, ensembles, process understanding, and multiple uncertainty sources, and future statements are conditional on forcing and scenario assumptions \cite{ipcc2021}. HCFCA therefore labels an output as a forecast only when it has a compatible held-out target, and labels a scenario as a conditional projection when the forcing path is an assumption rather than an observed future.

\subsection{Evaluation and calibration}
A predictive distribution $P$ should be evaluated with a proper scoring rule. Gneiting and Raftery show why strictly proper rules incentivize honest probability statements \cite{gneiting2007}. For a continuous target, the log score and continuous ranked probability score (CRPS) assess distributional quality; for a binary event, the Brier score evaluates probabilistic accuracy. Reliability, coverage, sharpness, and tail-event performance provide complementary information. No single metric should be allowed to conceal a failure in another.

The M4 competition found that forecast quality depends on series family, aggregation, horizon, and method, and that strong simple baselines remain difficult to beat \cite{makridakis2018}. DeepAR illustrates a probabilistic neural approach that produces distributions for related time series \cite{salinas2020}. The proposed benchmark inherits the lesson that an intelligent model must be compared with persistence, climatology, autoregression, state-space, and mechanistic baselines under the same split and information budget.

\subsection{The survey gap ledger}
The evidence supports seven gaps. G1 is a type gap: deterministic prediction, probabilistic forecasting, and scenario projection are frequently discussed together. G2 is a horizon gap: calendar time cannot be compared across domains without a characteristic timescale. G3 is an uncertainty gap: initial-state, observation, model, process, and shift uncertainty are confounded. G4 is a calibration gap: point skill is reported more often than probability quality. G5 is a transfer gap: AI gains shown in one domain may not survive a changed regime. G6 is an information gap: computation is not a substitute for missing observations. G7 is a decision gap: a skillful forecast may be poorly timed or miscalibrated for an operational decision.

The Survey decision is to study the boundary itself. Another leaderboard would report who won on one set of targets. HCFCA instead asks where and why skill survives, and whether uncertainty estimates track that survival.

\section{Research Questions and Planned Contributions}
The study will test three primary questions:
\begin{enumerate}
\item Does an AI or hybrid forecaster retain out-of-sample skill at longer normalized horizons than strong statistical and mechanistic baselines?
\item Can calibration and uncertainty perturbations identify an effective horizon at which a forecast ceases to add reliable information?
\item Does forecastability transfer across stable dynamical systems, weather and climate, economic series, and aggregate public-health series, or is it conditional on observations and regime stability?
\end{enumerate}

The planned contributions are fourfold. First, a typed forecast contract prevents an exact-state score from being compared with a scenario narrative. Second, normalized horizons make cross-domain curves interpretable without pretending that all domains share one clock. Third, a joint skill--calibration--shift atlas distinguishes better predictions from better uncertainty communication. Fourth, the information ladder and compute controls test whether an apparent AI gain comes from access to data, model inductive bias, or merely more computation.

\section{Idea Source Checkpoints and Direction Selection}
The Idea Agent considered four classes of alternatives. A chaos-only direction based on local divergence supplies a useful physical anchor, but cannot evaluate calibration or policy feedback. A scale-only direction compares parameter count and inference speed, but confuses computational capacity with information. A scenario-only direction is essential for conditional climate planning, but is not an exact event forecast. A general future oracle based on a language model has no stable target and is vulnerable to leakage and narrative overconfidence.

HCFCA combines the useful parts of the first three while rejecting the oracle framing. Its root mechanism is that predictive information decays through initial-state uncertainty, model approximation, observation noise, stochastic forcing, and distribution shift. An intelligent model can slow the empirical decay or characterize it more accurately, but cannot create an observation of an unobserved event. The primary idea is selected because it is novel enough to test a real boundary, aligned with the Survey evidence, feasible with public data, and falsifiable without assuming a positive result.

The idea evolved through four checkpoints. The untyped future became three output types. Calendar lead time became lead time divided by a characteristic domain timescale. A single accuracy score became a family of proper scores and calibration diagnostics. Finally, a model leaderboard became a shift-aware atlas with explicit information and resource controls.

\section{Problem Definition, Assumptions, and Hypotheses}
Let $d$ index domains, $m$ models, $t$ forecast origins, and $h$ lead times. The information set $X_{d,t}$ can include observations, metadata, and declared forcings, while $Y_{d,t+h}$ is the target. The benchmark assumes that each data source has an explicit version, timestamp, and release-vintage rule. It does not assume stationarity, perfect observation, or that one model class is best for all targets.

The main hypotheses are:
\begin{itemize}
\item H1: Dense observations and a model matched to the system extend useful skill, with the largest expected gain in observation-rich weather and selected orbital targets.
\item H2: Calibrated probabilistic ensembles retain decision-relevant value farther into the horizon than point forecasts, even when mean error grows.
\item H3: Initial-condition, observation, and distribution-shift perturbations shorten the effective horizon through distinguishable signatures.
\item H4: No model maintains exact-state superiority indefinitely; long-horizon outputs approach a climatological, stationary, or scenario-conditional baseline according to target type.
\end{itemize}

These are not claims about observed results. They are preregistered propositions whose failure would be informative.

\section{Study Design and Methods}
\subsection{Benchmark cohorts}
The benchmark contains five cohorts. The orbital cohort uses public ephemerides and satellite state vectors as a high-predictability control. It evaluates position and velocity errors under known dynamics and perturbed initial states. The weather cohort uses versioned reanalysis and archived forecast products, targeting temperature, pressure, wind, precipitation, and severe-event indicators. The climate cohort uses historical or hindcast segments for evaluation and treats future scenario runs as conditional projections. The economic cohort uses release-vintage-controlled M4 or FRED-style macroeconomic series. The public-health cohort uses aggregate epidemic forecasting-hub or surveillance series with reporting-delay metadata; it is methodological and does not generate individual risk predictions.

For target $j$ in domain $d$, estimate a characteristic timescale $\tau_{d,j}$ from autocorrelation decay, a seasonal period, or a preregistered dynamical estimate. Evaluate normalized horizons $h/\tau_{d,j}\in\{0.25,0.5,1,2,4,8\}$ when supported by the record length. Report calendar horizons as a secondary quantity. If two timescale estimators disagree materially, both are retained in sensitivity analysis rather than selected after seeing model performance.

\subsection{Rolling origins and information ladder}
The design uses rolling-origin temporal splits. Training and calibration precede each test window; an untouched final period is never used for model choice. Random shuffles are prohibited. Data revisions, sensor changes, and reporting delays are recorded. A leakage audit checks that a feature timestamp is earlier than the forecast origin and that scenario labels do not encode the held-out target.

Each origin is evaluated under five information conditions: full observations; masked observations; coarsened temporal or spatial observations; added observation noise; and delayed observations. For dynamical cohorts, initial states are also perturbed. The model perturbation condition varies parameterizations or stochastic ensemble members. Shift windows include normal temporal holdout, event-rich holdout, and a documented release, policy, or forcing transition when metadata permit.

\subsection{Forecasting systems and controls}
The minimum model set comprises persistence and seasonal climatology, AR/VAR or state-space models, a mechanistic ensemble where available, gradient-boosted lag features, a probabilistic autoregressive neural model, a transformer or graph model for high-dimensional fields, and a hybrid residual model. Every model sees the same information ladder and origin split. Parameter count, training-data volume, inference latency, and compute are recorded as covariates.

The hybrid model receives a mechanistic forecast and learns a residual distribution. This allows the study to test a meaningful role for AI: correcting model error or extracting observational structure while retaining physical constraints. The design does not require that the hybrid model win. A simple baseline that is stable, calibrated, and less expensive is a legitimate scientific outcome.

\subsection{Forecast scores}
For a predictive distribution $P_{d,m,t,h}$ and outcome $y_{t+h}$, use
\begin{equation}
 LS=-\log p(y_{t+h}),
\end{equation}
and
\begin{equation}
 CRPS(P,y)=\int_{-\infty}^{\infty}\left(F_P(z)-\mathbf{1}\{z\ge y\}\right)^2 dz.
\end{equation}
For a binary event, use $BS=(p-y)^2$. Continuous state errors use normalized RMSE and, for multivariate states, an energy score. Calibration is measured by reliability, probability integral transform diagnostics, interval coverage, and sharpness. For a nominal interval level $\alpha$, define the coverage error $CE=|\widehat{coverage}-\alpha|$.

Skill relative to the strongest eligible baseline is
\begin{equation}
 S_{d,m,h}=1-\frac{L_{d,m,h}}{L_{d,b^*,h}},
\end{equation}
where $b^*$ is chosen on validation data only. A model with positive point skill but poor calibration is not declared superior for a probabilistic decision.

\subsection{Forecastability and effective horizon}
Define a shift penalty $U_{d,m,h}$ from preregistered covariate drift and the degradation of calibration between validation and test windows. The descriptive atlas score is
\begin{equation}
 F_{d,m,h}=S_{d,m,h}(1-\widetilde{CE}_{d,m,h})(1-U_{d,m,h}),
\end{equation}
where calibration error is scaled using validation-defined bounds. The score is not proposed as a universal law; it is an auditable summary of a specific target and information set.

The effective horizon is
\begin{equation}
 H^*_{d,m}=\sup\{h:S_{d,m,h}>\delta,\;CE_{d,m,h}<\epsilon,\;\Delta F_{d,m,h}>0\}.
\end{equation}
Thresholds $\delta$ and $\epsilon$ are fixed before final testing. The positive-improvement criterion requires a block-bootstrap interval above zero, with multiplicity controlled over the target family. Block bootstrap over forecast origins preserves serial dependence.

As a model-independent companion, estimate the information curve
\begin{equation}
 I_d(h)=H(Y_{t+h})-H(Y_{t+h}\mid X_t).
\end{equation}
The information curve is not equated with model skill. A poor model can waste available information, while a miscalibrated model can appear sharp without being reliable.

\subsection{Hierarchical analysis}
Origin-level losses are modeled using
\begin{equation}
\begin{split}
 L_{d,m,h,w}={}&\beta_0+\beta_1(h/\tau_d)+\beta_2D_{shift}+\beta_3D_{obs}+\beta_4M_m\\
 &+\beta_5(h/\tau_d)D_{shift}+u_d+u_w+\varepsilon.
\end{split}
\end{equation}
Here $u_d$ and $u_w$ represent domain and rolling-window heterogeneity. The primary contrast is the difference in $H^*$ between HCFCA and the strongest baseline. Secondary contrasts test calibration, tail events, and compute-normalized performance. Results are reported per target and per domain before any pooled summary.

\subsection{Reproducibility and sample planning}
This proposal plans public-data evaluation rather than a human sample. A minimum release contains at least 20 independent rolling origins per target when the record permits, five information conditions, six normalized horizons, and two untouched evaluation windows. If a cohort cannot provide that support, its role is descriptive and it cannot dominate the pooled estimate. All data versions, splits, seeds, calibration maps, failed runs, and hardware metadata are registered. The exact origin count is frozen after data availability is audited and before final model comparison.

\section{Expected Outcome Branches and Conditional Conclusions}
Four outcomes are scientifically substantive. In Branch A, AI extends skill and calibration in observation-rich domains, but all forecast curves eventually decay. This would support a bounded extension of useful prediction. In Branch B, AI improves point skill while ensemble or conformal recalibration supplies most decision value; the result would favor uncertainty-aware deployment over a point-forecast leaderboard. In Branch C, gains are strong for stable orbital and medium-range weather targets but collapse under economic, public-health, or climate regime shifts; this would establish transferability as the limiting variable. In Branch D, strong statistical or mechanistic baselines match AI after leakage and compute controls; the positive contribution would be a sharper forecastability boundary and a demonstration that more complexity is not automatically more knowledge.

The strongest positive conclusion permitted by the design is that an intelligent machine extends a typed, horizon-bounded forecast under a declared information set. The strongest negative conclusion is not that the future is unknowable, but that a particular target and information set does not support reliable prediction beyond its measured effective horizon. Scenario projections remain useful conditional statements, not observations of a realized future.

\section{Additional Analysis Plan}
\subsection{Transfer, adaptation, and the meaning of a gain}
An AI model can obtain a favorable score by exploiting a stable correlation that is not a durable mechanism. HCFCA therefore separates three evaluation regimes. In interpolation, the test window resembles the training distribution. In temporal extrapolation, the test window is later and may contain new combinations of covariates. In structural shift, a documented measurement, policy, forcing, or reporting process changes. A model that performs well only in interpolation has a narrower claim than one that remains calibrated under extrapolation.

For each target, the protocol reports a transfer ratio,
\begin{equation}
 T_{d,m}=\frac{F_{d,m,h}^{shift}}{F_{d,m,h}^{stable}},
\end{equation}
at matched normalized horizons. The ratio is descriptive and is accompanied by uncertainty; it is not interpreted as a universal property of a model architecture. A low ratio can arise from a changed target mechanism, reduced observations, or calibration failure. These causes are separated using the information ladder and shift metadata.

Adaptation is treated as a new forecast problem, not as free information. A frozen model, a model recalibrated with an allowed recent window, and a model retrained with the post-shift window are reported separately. The adaptation window cannot contain the final evaluation outcomes. This design tests a practical question: whether an intelligent machine can recognize that its previous forecast distribution is obsolete and repair it using legitimately available information.

\subsection{Causal interpretation of information perturbations}
The perturbation conditions do not prove causal mechanisms by themselves, but they provide intervention-like tests of information dependence. If masking a predictor or delaying an observation changes $H^*$ while leaving the target and model class fixed, the forecast boundary is sensitive to that information channel. If adding noise changes the error curve but not calibration, uncertainty may be represented honestly even as skill declines. If calibration degrades only after a regime transition, shift rather than ordinary error growth is implicated.

The planned analysis compares observed perturbation effects with a null in which the removed information is redundant. Permutation and blockwise masking are performed within the training protocol and frozen before final evaluation. For the orbital and weather cohorts, physically meaningful perturbations are preferred over arbitrary feature deletion. For economic and public-health cohorts, release delays and reporting masks mimic information actually available at the origin. The result is a causal-structure test of the forecast pipeline, not a claim that an observational benchmark identifies every real-world mechanism.

\subsection{Rare events and asymmetric decisions}
Average loss can hide failure on the events for which a forecast is most valuable. Each cohort therefore defines a small number of target-specific tail events before testing, such as a high-percentile temperature, extreme precipitation, a large orbital deviation, a rapid epidemic increase, or an economic threshold crossing. Event definitions use training-era quantiles and cannot be tuned on the final period. Evaluation includes event Brier score, precision-recall behavior, coverage of prediction intervals, and conditional calibration by event prevalence.

When an action loss is available, define a threshold policy $a(p)$ from the predictive probability $p$ and evaluate expected loss on held-out outcomes. The threshold is selected from validation data only. This makes decision usefulness explicit: a cautious policy and an aggressive policy may prefer different calibrated distributions even when they have the same global log score. No action is executed by this proposal. The action-loss analysis only determines what a future operational deployment would need to measure.

\subsection{Uncertainty in the estimated horizon}
The horizon is a boundary estimated from finite, serially dependent origins. Confidence intervals are obtained by resampling contiguous origin blocks, with block length chosen from the largest relevant dependence estimate and fixed before final analysis. A second interval is obtained by varying the characteristic timescale estimator. The report presents the full curve together with $H^*$ rather than only a threshold point, because a shallow decline and a sharp collapse have different scientific meanings.

If the skill curve crosses the threshold more than once, the primary definition uses the last contiguous interval that satisfies the criteria, and the full non-monotone curve is retained. Non-monotonicity is expected in systems with seasonal windows or event-conditioned predictability. If no horizon satisfies the minimum skill criterion, $H^*$ is reported as below the first tested horizon; if all horizons satisfy it, the result is right-censored and not interpreted as infinite predictability.

\subsection{Power, multiplicity, and negative results}
The unit of replication is the forecast origin within a target, not an individual time point. Before the final run, a simulation-based power analysis uses the observed dependence structure of the training period and a range of effect sizes that are scientifically meaningful. The analysis is powered for the primary contrast in $H^*$ and for calibration degradation under shift, not for every model-target pair. False-discovery-rate control is applied across secondary target-specific hypotheses, and effect sizes are retained even when adjusted intervals include zero.

The proposal treats a negative result as a boundary measurement. If no AI model extends $H^*$ after matched information, compute, and leakage controls, that result narrows the claim that intelligence helps. If AI extends point skill but not calibrated skill, it identifies uncertainty representation as the bottleneck. If a baseline matches the learned model, the baseline becomes evidence about the target's accessible structure rather than a failed competitor. This prevents the final analysis from equating novelty with scientific success.

\subsection{Reporting protocol}
For each domain and target, publish: data version and provenance; forecast-origin diagram; information availability; model and parameter budget; normalized and calendar horizons; deterministic and probabilistic scores; calibration curves; shift penalty; bootstrap interval; compute and latency; and the exact rule used to determine $H^*$. A pooled summary is shown only after these panels. Every conclusion must name its target type, horizon, information condition, and shift regime. This reporting contract makes it possible to distinguish a genuine extension of predictive information from a better-looking but overconfident forecast.

\section{Risks, Limitations, and Human Review Requirements}
The principal risk is target leakage. Vintage controls, timestamp audits, and an untouched final period reduce but cannot eliminate it. A second risk is metric gaming: a model may optimize one proper score while failing tail calibration or decision utility. The protocol therefore reports several diagnostics and disallows a single composite score from replacing the underlying curves.

Timescale normalization can be misleading when a system has multiple attractors or abrupt regime changes. The proposal addresses this with target-specific timescales, alternative estimators, and shift-stratified results. The atlas score can also create false precision if its thresholds are chosen opportunistically; thresholds and bootstrap rules must be preregistered. Cross-domain pooling can conceal failure, so pooled estimates are secondary to domain and target panels.

Public-health data are affected by reporting delay, access, policy, and behavior. Outputs must remain aggregate methodological evaluations. No individual prediction, clinical decision, or public warning is generated by this package. Any operational use requires domain experts to examine data provenance, calibration, tail behavior, equity, and communication consequences. Climate scenarios must be labeled as conditional. Orbital accuracy must not be generalized to social or biological systems.

Human review is required before release of code or forecasts for high-consequence decisions. Reviewers should verify that every data source is legally accessible, every forecast origin precedes its target, every probability is calibrated on an appropriate reference period, and every claim is bounded by its target and horizon. The design is intended to make overclaiming difficult, not to replace scientific judgment.

\section{Conclusion}
The future is neither a single unknowable object nor a fully accessible sequence waiting for more computation. It contains tasks with very different predictability structures. Satellite trajectories can be propagated with high precision over useful horizons; weather systems require ensembles; climate information is often conditional on forcing scenarios; economic and public-health forecasts must confront feedback and regime change. An advanced intelligent machine can improve the representation of these possibilities and can extend practical skill when data and structure are available. It cannot remove missing information, chaotic amplification, or an unobserved intervention.

HCFCA turns that statement into a testable research program. It measures skill, calibration, information, shift, and compute along normalized horizons; compares AI with strong baselines; and reports the horizon where a forecast stops adding reliable information. The deliverable is therefore not a prophecy. It is a reproducible map of where prediction works, where it fails, and what kind of future statement the evidence can support.

\appendices
\section{Variables and Operational Definitions}
\begin{table*}[t]
\caption{Core variables and definitions}
\label{tab:variables}
\centering
\begin{tabular}{p{0.19\textwidth}p{0.70\textwidth}}
\toprule
Variable & Operational definition\\
\midrule
Target type & Deterministic state, probabilistic event/distribution, or scenario-conditional projection.\\
Characteristic time $\tau$ & Preregistered autocorrelation, seasonal, or dynamical timescale for each target.\\
Information quality & Full, masked, coarsened, noisy, delayed, or initial-state-perturbed information condition.\\
Skill $S$ & Relative improvement in held-out loss over the strongest validation-selected baseline.\\
Calibration error $CE$ & Absolute difference between empirical and nominal probability coverage, complemented by reliability.\\
Shift penalty $U$ & Frozen function of covariate drift and validation-to-test calibration degradation.\\
Atlas score $F$ & Descriptive product of baseline skill, scaled calibration quality, and shift robustness.\\
Effective horizon $H^*$ & Largest lead time satisfying skill, calibration, and positive-improvement criteria.\\
Decision usefulness & Secondary domain-specific evaluation against a declared action loss; never inferred from accuracy alone.\\
\bottomrule
\end{tabular}
\end{table*}

\section{Gap-to-Test Matrix}
\begin{table*}[t]
\caption{Accepted Survey gaps and corresponding tests}
\label{tab:gaps}
\centering
\begin{tabular}{p{0.11\textwidth}p{0.31\textwidth}p{0.45\textwidth}}
\toprule
Gap & Problem & Planned test\\
G1 & Forecast, prediction, and projection are conflated. & Typed output contract and compatible scoring sets.\\
G2 & Calendar horizon is not comparable across domains. & Target-specific $h/\tau$ curves and timescale sensitivity.\\
G3 & Uncertainty sources are confounded. & Information ladder, initial-state perturbation, model ensembles, and shift windows.\\
G4 & Accuracy is reported without probability quality. & Proper scores, coverage, reliability, sharpness, and tail-event metrics.\\
G5 & AI gains may not transfer. & Cross-domain and regime-shift evaluation with frozen calibration.\\
G6 & Compute is mistaken for information. & Parameter, data, latency, and compute covariates with matched baselines.\\
G7 & Statistical skill is not decision usefulness. & Declared action loss and expert review for operational targets.\\
\bottomrule
\end{tabular}
\end{table*}

\section{Evidence Coverage and Review Appendix}
The Survey evidence supports the following bounded claims. Lorenz supports sensitive dependence and nonperiodic deterministic flow \cite{lorenz1963,lorenz1969}. ECMWF supports the operational treatment of initial and model uncertainty with ensembles \cite{ecmwf}. Gneiting and Raftery support proper scoring rules for predictive distributions \cite{gneiting2007}. The IPCC supports multi-source uncertainty and conditional climate information \cite{ipcc2021}. GraphCast supports a concrete, bounded AI advance in medium-range global weather forecasting \cite{lam2023}. The M4 and DeepAR sources support strong baselines and probabilistic forecasting practice \cite{makridakis2018,salinas2020}. The Google Flu analysis supports caution about proxy drift and changing relationships \cite{lazer2014}.

These sources do not establish a universal forecastability scalar, an exact future event, or the superiority of the proposed HCFCA system. Those are design questions. The present package preserves an empty observed-results field and should not be read as an executed benchmark.

\begin{thebibliography}{00}
\bibitem{lorenz1963} E. N. Lorenz, ``Deterministic nonperiodic flow,'' \emph{J. Atmos. Sci.}, vol. 20, no. 2, pp. 130--141, 1963, doi: 10.1175/1520-0469(1963)020<0130:DNF>2.0.CO;2.
\bibitem{lorenz1969} E. N. Lorenz, ``The predictability of a flow which possesses many scales of motion,'' \emph{Tellus}, vol. 21, no. 3, pp. 289--307, 1969, doi: 10.1111/j.2153-3490.1969.tb00444.x.
\bibitem{ecmwf} European Centre for Medium-Range Weather Forecasts, ``Quantifying forecast uncertainty,'' accessed Sep. 1, 2026. [Online]. Available: https://www.ecmwf.int/en/research/modelling-and-prediction/quantifying-forecast-uncertainty
\bibitem{leutbecher2008} M. Leutbecher and T. N. Palmer, ``Ensemble forecasting,'' \emph{J. Comput. Phys.}, vol. 227, pp. 3515--3539, 2008, doi: 10.1016/j.jcp.2007.02.014.
\bibitem{gneiting2007} T. Gneiting and A. E. Raftery, ``Strictly proper scoring rules, prediction, and estimation,'' \emph{J. Amer. Stat. Assoc.}, vol. 102, no. 477, pp. 359--378, 2007, doi: 10.1198/016214506000001437.
\bibitem{ipcc2021} IPCC, ``Chapter 10: Linking global to regional climate change,'' in \emph{Climate Change 2021: The Physical Science Basis}, Cambridge Univ. Press, 2021. [Online]. Available: https://www.ipcc.ch/report/ar6/wg1/chapter/chapter-10/
\bibitem{lam2023} R. Lam \emph{et al.}, ``GraphCast: Learning skillful medium-range global weather forecasting,'' \emph{Nature}, vol. 619, pp. 533--538, 2023, doi: 10.1038/s41586-023-06185-3.
\bibitem{makridakis2018} S. Makridakis, E. Spiliotis, and V. Assimakopoulos, ``The M4 Competition: Results, findings, conclusion and way forward,'' \emph{Int. J. Forecasting}, vol. 34, no. 4, pp. 802--808, 2018, doi: 10.1016/j.ijforecast.2018.06.001.
\bibitem{salinas2020} D. Salinas, V. Flunkert, J. Gasthaus, and T. Januschowski, ``DeepAR: Probabilistic forecasting with autoregressive recurrent networks,'' \emph{Int. J. Forecasting}, vol. 36, no. 3, pp. 1181--1191, 2020, doi: 10.1016/j.ijforecast.2019.07.001.
\bibitem{hyndman2021} R. J. Hyndman and G. Athanasopoulos, \emph{Forecasting: Principles and Practice}, 3rd ed. Melbourne, Australia: OTexts, 2021.
\bibitem{lazer2014} D. Lazer, R. Kennedy, G. King, and A. Vespignani, ``The parable of Google Flu: Traps in big data analysis,'' \emph{Science}, vol. 343, no. 6176, pp. 1203--1205, 2014, doi: 10.1126/science.1248506.
\end{thebibliography}
\end{document}
